\documentclass[12pt]{article}
\usepackage[utf8]{inputenc}
\usepackage[T1]{fontenc}
\usepackage{hyperref}
\title{Wandji Language Resources for K-12 Curriculum Development}
\author{Compiled by Claude Sonnet 5 (claude-sonnet-5) for Andr\'as Kornai}
\date{}
\begin{document}
\maketitle

\section{Introduction}
Wandji (ISO 639-3 wdd; Glottocode wand1266) is also called Wanzi, Liwanzi or Bawandji. It is a Bantu language of Gabon (Guthrie B.501, Nzebi group), reported in Haut-Ogoou\'e and Ogoou\'e-Lolo. Wikipedia gives about 10,000 native speakers (2007). Joshua Project gives a larger figure of 22,000 for the people group, so the estimates differ. Searches found no school materials, primers or dictionaries in Wandji.

\section{Primary Educational Resources}
No primers, textbooks or school programmes for Wandji were found in search results.

\section{Linguistic and Reference Materials}
The only reference page found is an encyclopedic overview, which has basic classification and speaker figures only: \url{https://en.wikipedia.org/wiki/Wanzi_language}. No grammar, dictionary or orthography description was found.

\section{Religious and Community Resources}
Joshua Project lists Bible translation as started, with no audio Bible or Jesus Film: \url{https://joshuaproject.net/languages/wdd} and \url{https://joshuaproject.net/people_groups/15854/GB}. Its people-group text is not well sourced and should be treated with caution.

\section{Digital and Audio Resources}
Global Recordings Network lists evangelistic and Bible-teaching audio recordings: \url{https://globalrecordings.net/en/language/wdd}.

\section{Curriculum Development Context}
Very little published material appears to exist. A curriculum would need new orthography work, community consultation and fieldwork. The related Nzebi language may offer comparative material, but none was searched for here. French is the likely language of schooling in Gabon.
\end{document}